\documentclass[10pt,a4paper]{amsart}
\usepackage[margin=19mm]{geometry}
\usepackage{amsmath,amssymb,mathtools}
\usepackage[scr=rsfs,cal=euler]{mathalfa}
\usepackage{xfrac}
\usepackage[unicode,hidelinks,bookmarks=false]{hyperref}
\newcommand{\ie}{i.e.\ }
\newcommand{\cf}{cf.\ }
\newcommand{\resp}{respectively\ }
\newcommand{\Subsection}{\subsection}
\providecommand{\nobreakdash}{\nobreak\hbox{-}}
\setlength{\emergencystretch}{2em}
\multlinegap0pt
\usepackage{xcolor}
\usepackage{float}
\usepackage{enumerate}
\usepackage[normalem]{ulem}
\newtheorem{lemma}{Lemma}
\newcommand{\N}{\mathbb{N}}
\newcommand{\plimgG}[3]{\mathop {#1\text{-lim}}_{#2\in #3}\,}
\title{Multipliers and Disjointness from Mixing}
\author{Sohail Farhangi, Joel Moreira, Rigoberto Zelada}
\date{Source: arXiv:2604.12915v1 (CC BY 4.0)}
\begin{document}
\maketitle

\noindent\emph{Source and licence.} Multipliers and Disjointness from Mixing. Version arXiv:2604.12915v1 (CC BY 4.0), distributed by arXiv under the Creative Commons Attribution 4.0 International licence, http://creativecommons.org/licenses/by/4.0/, as recorded in the arXiv licence metadata for this exact version and shown on the abstract page https://arxiv.org/abs/2604.12915v1. Full licence notice: \texttt{licenses/arxiv-2604.12915-CC-BY-4.0.txt}.\par\smallskip

\noindent\textit{Editorial scope.} Counted unit: the article's lemma lem:IndependentJoiningsViaVdC (source0.tex lines 1614--1621). The article's complete original proof of it is delivered with it (source0.tex lines 1622--1678, one \texttt{proof} environment of 57 lines). No mathematics is written, repaired or re-derived: the statement and the proof are the article's own lines, in the article's own notation, and no retained line is substituted or re-flowed.  Retained uncounted as the source-local context the proof needs: lines 1570--1591.  Nothing was omitted from any retained span, so the package carries no omission record.  No retained line contains a citation, so the wrapper carries no bibliography and no bibliography entry.  No retained line contains a figure environment, an image-inclusion command or drawing code, so no image is delivered and the package declares no asset dependency.  Editorial formatting only, in addition: an AMS \texttt{amsart} preamble that restates the article's own declarations and macros used by the retained lines, namely the environments \texttt{lemma} on separate counters, together with the standard packages those lines need.  The retained environments print in this document as Lemma 1 in order of appearance, whereas the article prints the same items with the numbering of its own full document.  The complete native source of the article as distributed in the arXiv e-print for this exact version is bundled unchanged as \texttt{sources/198435/source0.tex}.
\begin{lemma} [Generalized van der Corput trick] \label{3.GeneralVdC}
 Let $S$ be a non-empty set, let $p_1,...,p_\ell\in\beta S$ be ultrafilters, let  $\mathcal H$ be a Hilbert space, and let $j\in\{1,...,\ell\}$. Let
 $$(x_{s_1,...,s_\ell})_{s_1,...,s_\ell\in S}$$
 be norm\text{-}bounded  sequences in $\mathcal H$.
If 
\begin{multline}\label{3.LimitIs0}
\plimgG{p_1}{s_1}{S}\cdots\plimgG{p_{j-1}}{s_{j-1}}{S}\plimgG{p_{j}}{r}{S}\plimgG{p_{j}}{t}{S}\plimgG{p_{j+1}}{s_{j+1}}{S}\cdots\plimgG{p_{\ell}}{s_{\ell}}{S}\\
\langle x_{s_1,...,s_{j-1},r,s_{j+1},...,s_\ell},x_{s_1,...,s_{j-1},t,s_{j+1},...,s_\ell}\rangle=0,
\end{multline}
then for any bounded sequence $(y_{s_1,...,s_{j-1},s_{j+1},...,s_\ell})_{s_1,...,s_{j-1},s_{j+1},...,s_\ell\in S}$ in $\mathcal H$, 
\begin{multline}\label{3."Weak"LimitIs0}
\plimgG{p_1}{s_1}{S}\cdots\plimgG{p_{j-1}}{s_{j-1}}{S}\plimgG{p_{j}}{s_{j}}{S}\plimgG{p_{j+1}}{s_{j+1}}{S}\cdots\plimgG{p_{\ell}}{s_{\ell}}{S}\\
\langle y_{s_1,...,s_{j-1},s_{j+1},...,s_\ell},x_{s_1,...,s_{j-1},s_{j},s_{j+1},...,s_\ell}\rangle=0.
\end{multline}
So, in particular,  \eqref{3."Weak"LimitIs0} implies that for any $y\in \mathcal H$,
$$\plimgG{p_1}{s_1}{S}\cdots\plimgG{p_{j-1}}{s_{j-1}}{S}\plimgG{p_{j}}{s_{j}}{S}\plimgG{p_{j+1}}{s_{j+1}}{S}\cdots\plimgG{p_{\ell}}{s_{\ell}}{S}\\
\langle y, x_{s_1,...,s_{j-1},s_{j},s_{j+1},...,s_\ell}\rangle =0,$$
and, hence, 
$$\plimgG{p_1}{s_1}{S}\cdots\plimgG{p_{j-1}}{s_{j-1}}{S}\plimgG{p_{j}}{s_{j}}{S}\plimgG{p_{j+1}}{s_{j+1}}{S}\cdots\plimgG{p_{\ell}}{s_{\ell}}{S}
 x_{s_1,...,s_{j-1},s_{j},s_{j+1},...,s_\ell}=0$$
 weakly. 
\end{lemma}

\begin{lemma}\label{lem:IndependentJoiningsViaVdC}
Let $(G,+)$ be a countable abelian group, let $\mathcal X$, $\mathcal Y$, and $\mathcal Z$ be measure-preserving systems, and let $\mathcal U\subseteq G^*$ be a non-empty family of  ultrafilters. Suppose that $\mathcal Z$ is $\mathcal U$-mixing and that $\mathcal Y\perp \mathcal Z$. Let $\sigma$ be a joining of $\mu$ and $\nu$.
Then, for any $p_1,...,p_\ell\in \mathcal U$,  any $f_1,...,f_\ell\in L^\infty(X,\mu)$, any $F\in L^\infty(Y,\nu)$,   any $H\in L^{\infty}(Z,\rho)$, and any joining of the form $\lambda=\sigma\lor \rho$, one has 
\begin{multline}\label{eq:IndependentJoiningVdC}
\int_{X\times Y\times Z}H(z)F(y)\prod_{j=1}^\ell T^{p_j}f_j(x)\text{d}\lambda(x,y,z)\\
=\int_{X\times Y} F(y) \prod_{j=1}^\ell T^{p_j}f_j(x)\text{d}\sigma(x,y)\int_Z H(z)\text{d}\rho(z).
\end{multline}
\end{lemma}

\begin{proof}
It suffices to show that when $\int_z H\text{d}\rho=0$, the left-hand side of \eqref{eq:IndependentJoiningVdC} equals zero. 
To do this, we proceed by induction on $\ell$.\\

$\bullet$ \textit{ Case $\ell=1$.} When $\ell=1$, we have 
\begin{multline}\label{eq:VdCK=1.1}
\int_{X\times Y\times Z} T^{p_1}f_1(x)F(y)H(z)\text{d}\lambda(x,y,z)=
\plimgG{p_1}{g_1}{G}\int_{X\times Y\times Z} T^{g_1}f_1(x)F(y)H(z)\text{d}\lambda(x,y,z) \\
=\plimgG{p_1}{g_1}{G}\int_{X\times Y\times Z} f_1(x)S^{-g_1}F(y)R^{-g_1}H(z)\text{d}\lambda(x,y,z)
\end{multline}
By Lemma \ref{3.GeneralVdC} applied with $\ell=2$ and $p_2=\{A\subseteq G\,|\,0\in G\}$ (so, in particular, $p_1=p_1+p_2$), the expression in \eqref{eq:VdCK=1.1} equals zero provided that 
\begin{multline}\label{eq:VdCk=1.2}
0=\plimgG{p_1}{r}{G}\plimgG{p_1}{t}{G}\plimgG{p_2}{g}{G}\int_{X\times Y\times Z} S^{-r-g}F(y)R^{-r-g}H(z)S^{-t-g}F(y)R^{-t-g}H(z)\text{d}\lambda(x,y,z)\\
=\plimgG{p_1}{r}{G}\plimgG{p_1}{t}{G}\int_{X\times Y\times Z} S^{-r}F(y)R^{-r}H(z)S^{-t}F(y)R^{-t}H(z)\text{d}\lambda(x,y,z)=0.
\end{multline}
Since $\mathcal{Y}\perp\mathcal{Z}$, we see that $\lambda$ projects to $\nu\otimes\rho$ on $Y\times Z$.
Since $\mathcal Z$ is $\mathcal U$-mixing, we obtain
\begin{multline}
    \plimgG{p_1}{r}{G}\plimgG{p_1}{t}{G}\int_{X\times Y\times Z} S^{-r}F(y)R^{-r}H(z)S^{-t}F(y)R^{-t}H(z)\text{d}\lambda(x,y,z)\\
    = \plimgG{p_1}{r}{G}\plimgG{p_1}{t}{G}\left(\int_{Y}FS^{-t+r}F\text{d}\nu\int_Z HR^{-t+r}H\text{d}\rho\right)\\
    =\left(\plimgG{p_1}{r}{G}\plimgG{p_1}{t}{G}\int_{Y}FS^{-t+r}F\text{d}\nu\right)\left(\int_X H\text{d}\rho\right)^2=0.
\end{multline}
Thus, \eqref{eq:IndependentJoiningVdC} holds when $\ell=1$.\\ 

$\bullet$ \textit{ Case $\ell=k+1$.} Let $k\in\N$ and suppose now that \eqref{eq:IndependentJoiningVdC} holds for any $\tilde f_1,...,\tilde f_k\in L^\infty(X,\mu)$, any $\tilde F\in L^\infty(Y,\nu)$, and any $\tilde H\in L^\infty(Z,\rho)$. 
Set $\ell=k+1$ and  take $f_1,...,f_\ell\in L^\infty(X,\mu)$, $F\in L^\infty(Y,\nu)$, and $H\in L^\infty(Z,\rho)$ such that $\int_X H\text{d}\rho=0$.
Note that, by Theorem \ref{3.GeneralVdC} and the $T\times S\times R$ invariance of $\lambda$, 
\begin{multline*}
0
=  \int_{X\times Y\times Z} \prod_{j=1}^\ell T^{p_j}f_j(x)(F(y)H(z))\,\text{d}\lambda(x,y,z)\\
= \plimgG{p_\ell}{g_\ell}{G}\cdots \plimgG{p_1}{g_1}{G} \int_{X\times Y\times Z}\prod_{j=1}^\ell T^{-\sum_{i\neq j}g_i}f_j(S^{-\sum_{i=1}^\ell g_i}FR^{-\sum_{i=1}^\ell g_i}H)\,\mathrm{d}\rho,
\end{multline*}
provided that
\begin{multline*}
0
= \plimgG{p_\ell}{r_\ell}{G}\plimgG{p_\ell}{t_\ell}{G} \plimgG{p_k}{g_k}{G}\cdots\plimgG{p_1}{g_1}{G} \\
\int_{X\times Y\times Z} R^{-\sum_{i=1}^k g_i}\Big(HR^{-t_\ell+r_\ell}H\Big)S^{-\sum_{i=1}^k g_i}\Big( FS^{-t_\ell+r_\ell}F\Big)
\left(\prod_{j=1}^kT^{-\sum_{\substack{i=1\\ i\neq j}}^k g_i}(f_jT^{-t_\ell+r_\ell}f_j)\right)\text{d}\lambda.
\end{multline*} 
For each $g\in G$, Set
\[
\tilde H_g:=HR^{-g}H,
\;\tilde F_g := F S^{-g}F\;,
\qquad
\tilde f_{j,g} := f_jT^{-g}f_j,\quad j=1,\dots,k.
\]
By the inductive hypothesis and since $\mathcal Z$ is $\mathcal U$-mixing, we obtain
\begin{multline*}
\plimgG{p_\ell}{r_\ell}{G}\plimgG{p_\ell}{t_\ell}{G}\plimgG{p_k}{g_k}{G}\cdots\plimgG{p_1}{g_1}{G}\\
\int_{X\times Y\times Z} R^{-\sum_{i=1}^k g_i}\tilde H_{t_\ell-r_\ell}S^{-\sum_{i=1}^kg_i}\tilde F_{t_\ell-r_\ell}\left(\prod_{j=1}^kT^{-\sum_{\substack{i=1\\ i\neq j}}^k g_i}\tilde f_{j,(t_\ell-r_\ell)}\right)\text{d}\lambda\\
= \plimgG{p_\ell}{r_\ell}{G}\plimgG{p_\ell}{t_\ell}{G}
\int_{X\times Y\times Z} \tilde H_{t_\ell-r_\ell}(z)\tilde F_{t_\ell-r_\ell}(y)\prod_{j=1}^kT^{p_j}\tilde f_{j,(t_\ell-r_\ell)}(x)\text{d}\lambda(x,y,z)\\
=\plimgG{p_\ell}{r_\ell}{G}\plimgG{p_\ell}{t_\ell}{G}
\left(\int_{Z} HT^{-t_\ell+r_\ell}H\text{d}\rho(z)\right)\left(\int_{X\times Y}\tilde F_{t_\ell-r_\ell}(y)\prod_{j=1}^kT^{p_j}\tilde f_{j,(t_\ell-r_\ell)}(x)\text{d}\sigma(x,y)\right)=0,
\end{multline*}
which completes the induction.
\end{proof}

\end{document}
